%============================================================
% CH2 · PATTERN 3 — EXERCISES, HINTS, ANSWERS
%============================================================
\subsection{Pattern 3: Exercises and Going Further}
\label{ssec:ch2-p3-ex}

\begin{patternbox}{Pattern 3 Summary Card}
\renewcommand{\arraystretch}{1.25}
\begin{tabularx}{\linewidth}{@{}>{\bfseries}l X@{}}
Recognise & many queries about ``below $u$'', ``the subtree of $u$'', ``the $k$-th visited from $u$''; a peeling order; vertices numbered so that subtrees are consecutive \\
Record & one iterative DFS: $\tin{}$/$\tout{}$ (one clock), $\mathrm{pos}$/$\mathrm{sz}$/\texttt{order} (pre-order), \texttt{post} (finishing order), depth \\
Ancestor & $\tin{u} \le \tin{v} \wedge \tout{v} \le \tout{u}$ --- also ``which child of $y$ leads to $x$'' \\
Subtree & block \texttt{order[pos[u] .. pos[u]+sz[u]-1]}; with groups: sorted entry times per group + binary search, or the offline counter difference \\
Orders & pre-order / BFS prefixes are connected; post-order prefixes are unions of subtrees; reversed post-order of a DAG is topological \\
Top traps & \texttt{tin} used as an array index; unsorted children; $U$ itself forgotten; recursion at depth $2\cdot10^5$; products of sizes in \texttt{int} \\
\end{tabularx}
\end{patternbox}

\exlegend
\begin{enumerate}[leftmargin=*, label=\textbf{E3.\arabic*}, itemsep=6pt]
    \item \exeasy\ \pconcept \exown{Ancestor queries.} A tree with $n$ vertices is given by its $n-1$ edges and rooted at 1. Answer $q$ queries ``$u$ $v$'': does $u$ lie on the path from the root to $v$ ($u = v$ counts)? ($n, q \le 2\cdot10^5$.)
    \item \exeasy\ \pclassic \exprob{Takahashi Tour}{AtCoder ABC213 D}{https://atcoder.jp/contests/abc213/tasks/abc213_d}. A tree of $N$ cities. Starting in city 1, Takahashi always moves to the unvisited neighbour with the smallest number; if there is none, he goes back to the city he came from (and stops when that is impossible in city 1). Print every city he is in, in order, including city 1 at both ends. ($N \le 2\cdot10^5$.)
    \item \exmedium\ \pconcept \exown{Edge types after the fact.} Run an iterative DFS on a directed graph and store only $\tin{}$, $\tout{}$ and, for every vertex, the id of the edge through which it was entered. Afterwards classify every edge as tree, back, forward or cross --- without colours --- and verify the result against the colour method of Lesson~\ref{ssec:ch2-dfs-tree} on random graphs with self-loops and parallel edges.
    \item \exmedium\ \ptricky \cf{1600}\ \exprob{Kuro and Walking Route}{Codeforces 979C}{https://codeforces.com/problemset/problem/979/C}. A tree with $n$ towns and two special towns $x \ne y$. Count the ordered pairs $(u, v)$, $u \ne v$, whose path does \emph{not} visit $x$ and later $y$. ($n \le 3\cdot10^5$.)
    \item \exmedium\ \precent \cf{1800}\ \exprob{Iris and the Tree}{Codeforces 2006B}{https://codeforces.com/problemset/problem/2006/B}. A rooted tree whose vertices are numbered in DFS order (every subtree is a range of consecutive numbers); edge $(i, p_i)$ has an unknown non-negative weight $t_i$, and $\sum t_i = w$. After each of $n-1$ events ``$t_x = y$'' print $\sum_{i=1}^{n} \max \operatorname{dist}(i, i \bmod n + 1)$, where each maximum is taken independently over the still unknown weights. ($\sum n \le 2\cdot10^5$, $w \le 10^{12}$.)
\end{enumerate}

\subsubsection*{Hints}
\begin{enumerate}[leftmargin=*, label=\textbf{E3.\arabic*}]
    \item ``On the path from the root to $v$'' is ``ancestor of $v$''. One DFS, two comparisons per query.
    \item The walk prints a vertex when it is entered and once more each time a child's subtree is finished. In the iterative DFS these are exactly the push and the moment the top is popped and the parent becomes the top again.
    \item A tree edge is the edge recorded for its head. For every other edge $u \to v$ the intervals of $u$ and $v$ are nested one way, the other way, or disjoint.
    \item Root the tree at $y$. A bad path starts in the subtree of $x$ and ends outside the branch of $y$ that contains $x$. Which child of $y$ is that branch? The answer can exceed $2^{31}$.
    \item The subtree of $x$ is the range $x..x+\mathrm{sz}[x]-1$. A path $i \to i+1$ uses the edge above $x$ iff it enters or leaves that range. How many paths do that? Then think about where to put the unknown weight.
\end{enumerate}

\subsubsection*{Answers}

\begin{solutionbox}[E3.1 Ancestor queries]
$u$ is on the root path of $v$ iff $u$ is an ancestor of $v$ (or $v$ itself): the parenthesis test. The DFS is iterative because the tree may be a path; the parent is skipped since the input is an undirected edge list.
\begin{lstlisting}[style=cpp]
#include <bits/stdc++.h>
using namespace std;

// E3.1: tree rooted at 1, q queries "u v": does u lie on the path from the root to v?
// One iterative DFS (one clock), then two comparisons per query.
int main() {
    ios::sync_with_stdio(false);
    cin.tie(nullptr);
    int n;
    cin >> n;
    vector<vector<int>> adj(n + 1);
    for (int i = 0; i < n - 1; i++) {
        int a, b;
        cin >> a >> b;
        adj[a].push_back(b);
        adj[b].push_back(a);
    }
    vector<int> tin(n + 1), tout(n + 1), par(n + 1, 0), st(1, 1);
    vector<size_t> it(n + 1, 0);
    int timer = 0;
    tin[1] = timer++;
    while (!st.empty()) {
        int u = st.back();
        if (it[u] < adj[u].size()) {
            int v = adj[u][it[u]++];
            if (v == par[u]) continue;
            par[v] = u; tin[v] = timer++;
            st.push_back(v);
        } else {
            tout[u] = timer++;
            st.pop_back();
        }
    }
    int q;
    cin >> q;
    while (q--) {
        int u, v;
        cin >> u >> v;
        cout << (tin[u] <= tin[v] && tout[v] <= tout[u] ? "YES" : "NO") << "\n";
    }
}
\end{lstlisting}
On the tree $T$ of the pattern theory with the queries $(3,9)$, $(2,6)$, $(6,6)$, $(1,7)$, $(8,6)$ it prints \textbf{YES NO YES YES NO}. It agrees with a brute force (BFS parents from the root, climb from $v$ looking for $u$) on $500$ random trees with $n \le 10$ and random labels.
\end{solutionbox}

\begin{solutionbox}[E3.2 Takahashi Tour]
Sort the neighbour lists. The tour is the \textbf{Euler walk} of the DFS tree: a vertex is printed on entry and again after each of its children returns, $2N - 1$ numbers in total. With the iterator-style iterative DFS, ``entering $v$'' is the push, and ``coming back'' is the pop after which the new top is the parent.
\begin{lstlisting}[style=cpp]
#include <bits/stdc++.h>
using namespace std;

// E3.2 Takahashi Tour (ABC213 D): print the vertex when it is entered and again
// every time the walk comes back to it from a child. Neighbours sorted; iterative DFS.
int main() {
    ios::sync_with_stdio(false);
    cin.tie(nullptr);
    int n;
    cin >> n;
    vector<vector<int>> adj(n + 1);
    for (int i = 0; i < n - 1; i++) {
        int a, b;
        cin >> a >> b;
        adj[a].push_back(b);
        adj[b].push_back(a);
    }
    for (auto& l : adj) sort(l.begin(), l.end());
    vector<int> par(n + 1, 0), st(1, 1), walk(1, 1);
    vector<size_t> it(n + 1, 0);
    while (!st.empty()) {
        int u = st.back();
        if (it[u] < adj[u].size()) {
            int v = adj[u][it[u]++];
            if (v == par[u]) continue;
            par[v] = u;
            st.push_back(v);
            walk.push_back(v);                     // entering v
        } else {
            st.pop_back();
            if (!st.empty()) walk.push_back(st.back());   // back in the parent
        }
    }
    for (size_t i = 0; i < walk.size(); i++) cout << walk[i] << (i + 1 < walk.size() ? ' ' : '\n');
}
\end{lstlisting}
Both official samples are reproduced: \texttt{1 2 4 2 1 3 1} and \texttt{1 2 1 3 1 4 1 5 1}. A literal simulation of the statement (scan all neighbours at every step) gives the same walk on $500$ random trees with $N \le 10$; the path of $2\cdot10^5$ cities (walk of $399{,}999$ numbers) runs without stack trouble.
\end{solutionbox}

\begin{solutionbox}[E3.3 Edge types after the fact]
After the DFS, for an edge $i = (u \to v)$:
\begin{center}\small
\begin{tabular}{@{}ll@{}}
\toprule
\textbf{test (in this order)} & \textbf{type} \\
\midrule
$v$ was entered through edge $i$ & tree \\
$\tin{v} \le \tin{u}$ and $\tout{u} \le \tout{v}$ ($v$ ancestor of $u$, or $u = v$) & back \\
$\tin{u} < \tin{v}$ and $\tout{v} < \tout{u}$ ($v$ a proper descendant) & forward \\
otherwise (disjoint intervals; then $\tout{v} < \tin{u}$) & cross \\
\bottomrule
\end{tabular}
\end{center}
\emph{Why it matches the colour method.} When the colour method examines $u \to v$, $u$ is gray. A gray $v$ is an ancestor (back); a black $v$ with later entry is a descendant (forward); a black $v$ with earlier entry is finished and unrelated (cross); a white $v$ becomes a child --- but a \emph{second} parallel edge to the same child, examined later, finds it black and descendant: forward. That is why ``tree'' must be decided by the edge id, not by \texttt{par[v] == u}.
\begin{lstlisting}[style=cpp]
#include <bits/stdc++.h>
using namespace std;

// E3.3: classify the edges of a directed graph AFTER the DFS, from tin/tout and the tree-edge ids
// only, and compare with the colour method of Lesson "The DFS Tree" on random graphs.
// Edge i = (a, b). Types: 0 tree, 1 back, 2 forward, 3 cross.
typedef vector<pair<int, int>> Edges;

// one iterative DFS (outer loop over 1..n, edges in input order); returns tin, tout, treeEdge[v]
void stamps(int n, const Edges& e, vector<int>& tin, vector<int>& tout, vector<int>& treeEdge) {
    vector<vector<int>> out(n + 1);                     // edge ids leaving each vertex
    for (int i = 0; i < (int)e.size(); i++) out[e[i].first].push_back(i);
    tin.assign(n + 1, -1); tout.assign(n + 1, -1); treeEdge.assign(n + 1, -1);
    vector<size_t> it(n + 1, 0);
    int timer = 0;
    for (int s = 1; s <= n; s++) {
        if (tin[s] >= 0) continue;
        tin[s] = timer++;
        vector<int> st(1, s);
        while (!st.empty()) {
            int u = st.back();
            if (it[u] < out[u].size()) {
                int id = out[u][it[u]++], v = e[id].second;
                if (tin[v] < 0) { tin[v] = timer++; treeEdge[v] = id; st.push_back(v); }
            } else { tout[u] = timer++; st.pop_back(); }
        }
    }
}

vector<int> classifyByTimes(int n, const Edges& e) {
    vector<int> tin, tout, treeEdge, type(e.size());
    stamps(n, e, tin, tout, treeEdge);
    for (int i = 0; i < (int)e.size(); i++) {
        int u = e[i].first, v = e[i].second;
        if (treeEdge[v] == i) type[i] = 0;                                  // v was entered through i
        else if (tin[v] <= tin[u] && tout[u] <= tout[v]) type[i] = 1;      // v ancestor of u (or u = v)
        else if (tin[u] < tin[v] && tout[v] < tout[u]) type[i] = 2;        // v descendant of u
        else type[i] = 3;                                                   // disjoint intervals
    }
    return type;
}

// reference: the colour method, classifying each edge at the moment it is examined (recursive, small n)
vector<vector<int>> outR;
vector<int> colR, tinR, typeR;
Edges eR;
int timerR;
void dfsR(int u) {
    colR[u] = 1; tinR[u] = timerR++;
    for (int id : outR[u]) {
        int v = eR[id].second;
        if (colR[v] == 0) { typeR[id] = 0; dfsR(v); }
        else if (colR[v] == 1) typeR[id] = 1;
        else typeR[id] = (tinR[u] < tinR[v]) ? 2 : 3;
    }
    colR[u] = 2; timerR++;
}
vector<int> classifyByColours(int n, const Edges& e) {
    eR = e; outR.assign(n + 1, vector<int>()); colR.assign(n + 1, 0); tinR.assign(n + 1, -1);
    typeR.assign(e.size(), -1); timerR = 0;
    for (int i = 0; i < (int)e.size(); i++) outR[e[i].first].push_back(i);
    for (int s = 1; s <= n; s++) if (!colR[s]) dfsR(s);
    return typeR;
}

int main() {
    const char* name[4] = {"tree", "back", "forward", "cross"};
    // the graph D of Part A, edges in adjacency order
    Edges D = {{1, 2}, {1, 5}, {1, 6}, {2, 3}, {2, 4}, {3, 1}, {4, 3}, {5, 4}, {5, 6}, {6, 7}, {7, 6}, {8, 7}};
    vector<int> t = classifyByTimes(8, D);
    for (int i = 0; i < (int)D.size(); i++) cout << D[i].first << "->" << D[i].second << " " << name[t[i]] << "\n";
    // random check, self-loops and parallel edges included
    mt19937 rng(2024);
    int bad = 0;
    for (int test = 0; test < 20000; test++) {
        int n = rng() % 7 + 1, m = rng() % 13;
        Edges e;
        for (int i = 0; i < m; i++) e.push_back({(int)(rng() % n) + 1, (int)(rng() % n) + 1});
        if (classifyByTimes(n, e) != classifyByColours(n, e)) bad++;
    }
    cout << "random graphs: 20000, mismatches: " << bad << "\n";
}
\end{lstlisting}
\begin{lstlisting}[style=trace]
output: 1->2 tree      1->5 tree      1->6 forward   2->3 tree
        2->4 tree      3->1 back      4->3 cross     5->4 cross
        5->6 tree      6->7 tree      7->6 back      8->7 cross
        random graphs: 20000, mismatches: 0          (one edge per line in the real output)
\end{lstlisting}
The classification of $D$ is exactly the one of Lesson~\ref{ssec:ch2-dfs-tree}, and the $20{,}000$ random graphs ($n \le 7$, up to 12 edges, loops and parallel edges allowed) produce no mismatch.
\end{solutionbox}

\begin{solutionbox}[E3.4 Kuro and Walking Route]
\textbf{Model.} Root the tree at $y$. The path $u \to v$ visits $x$ and \emph{then} $y$ iff $u$ lies in the subtree of $x$ (then the path from $u$ climbs through $x$ towards $y$) and $v$ lies outside the branch of $y$ that contains $x$ (so the path must climb all the way up to $y$). That branch is the child $c$ of $y$ that is an ancestor of $x$ --- one ancestor test per child. Bad pairs: $\mathrm{sz}[x] \cdot (n - \mathrm{sz}[c])$; answer $n(n-1) - \mathrm{sz}[x](n - \mathrm{sz}[c])$.

\textbf{The trap.} $n(n-1)$ reaches $9\cdot10^{10}$: every factor must be \texttt{long long} (the star with centre $x$ gives $(n-1)^2 = 89{,}999{,}400{,}001$ for $n = 3\cdot10^5$).
\begin{lstlisting}[style=cpp]
#include <bits/stdc++.h>
using namespace std;

// E3.4 Kuro and Walking Route (CF 979C). Root the tree at y. A pair (u, v) is bad iff the path
// u -> v meets x and then y: u lies in x's subtree, v lies outside the branch of y that contains x.
// That branch is the child c of y with c an ancestor of x (tin/tout test).
// Answer = n(n-1) - sz[x] * (n - sz[c]), in 64-bit.
int main() {
    ios::sync_with_stdio(false);
    cin.tie(nullptr);
    int n, x, y;
    cin >> n >> x >> y;
    vector<vector<int>> adj(n + 1);
    for (int i = 0; i < n - 1; i++) {
        int a, b;
        cin >> a >> b;
        adj[a].push_back(b);
        adj[b].push_back(a);
    }
    vector<int> tin(n + 1), tout(n + 1), sz(n + 1, 1), par(n + 1, 0), st(1, y);
    vector<size_t> it(n + 1, 0);
    int timer = 0;
    tin[y] = timer++;
    while (!st.empty()) {                       // iterative: n = 3e5 and the tree may be a path
        int u = st.back();
        if (it[u] < adj[u].size()) {
            int v = adj[u][it[u]++];
            if (v == par[u]) continue;
            par[v] = u; tin[v] = timer++;
            st.push_back(v);
        } else {
            tout[u] = timer++;
            if (par[u]) sz[par[u]] += sz[u];
            st.pop_back();
        }
    }
    int c = -1;
    for (int v : adj[y])
        if (tin[v] <= tin[x] && tout[x] <= tout[v]) c = v;   // the branch of y holding x
    long long N = n;
    cout << N * (N - 1) - (long long)sz[x] * (N - sz[c]) << "\n";
}
\end{lstlisting}
Official samples: \textbf{5} and \textbf{4}. A brute force over all ordered pairs (BFS path, check the positions of $x$ and $y$) agrees on $500$ random trees with $n \le 9$. Large checks: the path $1 - 2 - \dots - 3\cdot10^5$ with $x = 1$, $y = n$ gives $89{,}999{,}699{,}999 = n(n-1) - 1$, the star gives the value above.
\end{solutionbox}

\begin{solutionbox}[E3.5 Iris and the Tree]
\textbf{The block.} The numbering is a pre-order, so the subtree of $x$ is the range $B_x = [x, x + \mathrm{sz}[x] - 1]$, with sizes from the loop $i = n..2$ ($p_i < i$). The edge $e_x = (x, p_x)$ lies on the path between $a$ and $b$ iff exactly one of them is in $B_x$. The $n$ paths $i \to i \bmod n + 1$ walk once around the numbers $1, 2, \dots, n, 1$; they cross the boundary of the range $B_x$ exactly twice: path $x - 1$ (entering) and path $x + \mathrm{sz}[x] - 1$ (leaving; if that is $n$, the path $n \to 1$, and $1 \notin B_x$ since $x \ge 2$). So \textbf{every edge lies on exactly two paths}.

\textbf{The maximum.} Let $K$ be the sum of the known weights. A path that still contains an unknown edge can receive the whole rest $w - K$ on it; a path without unknown edges has a fixed length. Summing, the known weights contribute $2K$, and every ``open'' path contributes $w - K$:
\[
    \text{answer} = 2K + (\#\text{open paths}) \cdot (w - K).
\]
Keep, per path, the number of unknown edges on it; an event decrements the two counters of $e_x$, and a counter reaching 0 closes a path. \bigO{n} per test case.
\begin{lstlisting}[style=cpp]
#include <bits/stdc++.h>
using namespace std;

// E3.5 Iris and the Tree (CF 2006B). Vertices are numbered in DFS pre-order, so the subtree of x
// is the block x .. x+sz[x]-1. The edge (x, p_x) lies on exactly two of the n paths i -> i mod n + 1:
// path x-1 (enters the block) and path x+sz[x]-1 (leaves it; index n means n -> 1).
// Answer = 2 * known + (#paths that still have an unknown edge) * (w - known).
int main() {
    ios::sync_with_stdio(false);
    cin.tie(nullptr);
    int T;
    cin >> T;
    while (T--) {
        int n;
        long long w;
        cin >> n >> w;
        vector<int> p(n + 1, 0), sz(n + 1, 1), unknown(n + 1, 0);
        for (int i = 2; i <= n; i++) cin >> p[i];
        for (int i = n; i >= 2; i--) sz[p[i]] += sz[i];          // p_i < i
        for (int x = 2; x <= n; x++) { unknown[x - 1]++; unknown[x + sz[x] - 1]++; }
        int open = 0;                                              // paths with an unknown edge
        for (int i = 1; i <= n; i++) open += (unknown[i] > 0);
        long long known = 0;
        for (int e = 0; e < n - 1; e++) {
            int x;
            long long y;
            cin >> x >> y;
            known += y;
            if (--unknown[x - 1] == 0) open--;
            if (--unknown[x + sz[x] - 1] == 0) open--;
            cout << 2 * known + (long long)open * (w - known) << (e + 2 < n ? ' ' : '\n');
        }
    }
}
\end{lstlisting}
The official sample is reproduced: \texttt{2000000000000}, \texttt{25 18 18}, \texttt{449 302 247 200 200}, \texttt{4585 4473 2681 1567 1454 1322 1094 1022 1022}. A brute force that, after every event, walks each of the $n$ paths by climbing parents and adds ``known on the path $+ (w - K)$ if an unknown edge is on it'' agrees on $500$ random inputs (up to 3 test cases, $n \le 9$, random DFS-numbered trees). On a path of $2\cdot10^5$ vertices with $w = 10^{12}$ the first answer is $1.99999\cdot10^{17}$ --- \texttt{long long} is required.
\end{solutionbox}
